\section{Sensitivity analyses}
\label{si:sensitivities}

Three choices fixed by the registered design could, in principle, move the bound on the advantage of the electronic
descriptors, $A = \mathrm{MAE}_\mathrm{S} - \mathrm{MAE}_\mathrm{D}$, at the largest training set: the inclusion of
vacancies that leave a magnetic moment on the cell, the cap on training length and the choice of architecture
(Figure~\ref{fig:sens}; \texttt{docs/a5\_notes.md}). All three were run for the originally registered network
(v1) only; no v2 sensitivity runs exist, so the v2 statement rests on its primary result alone (an upper bound of
\numSensCappedPostHocVTwoPrimarySixHundredFiftyFourPrimaryAUpperOneSided{}~eV at \numCfgBudgetsHostsItemSix{}
hosts, post hoc).

The database treats every host as non-magnetic, but a few neutral vacancies leave a moment of more than
\numSensHighMomentVOneThresholdMuB{}~$\mu_\mathrm{B}$ on the cell, as two electrons localised on reducible cations
would; such polaron-like configurations are the ones that a non-magnetic target describes least well. Retraining S
and D with these entries removed from training and test (extension) changes $A$ on the common test sites from
\numSensHighMomentVOnePrimaryCommonSitesAMean{} to \numSensHighMomentVOneHighMomentAMean{}~eV, a paired change of
\numSensHighMomentVOneChangeHighMomentMinusPrimaryAMean{}~eV
\ci{\numSensHighMomentVOneChangeHighMomentMinusPrimaryACiLo{}}{\numSensHighMomentVOneChangeHighMomentMinusPrimaryACiHi{}},
and the one-sided upper bound moves from \numSensHighMomentVOnePrimaryCommonSitesAUpperOneSided{} to
\numSensHighMomentVOneHighMomentAUpperOneSided{}~eV, below the tolerance in both cases. The excluded entries are
rare in any test set, so the exclusion cannot move the pooled bound by much, and it does not.

Every primary run that stopped at the cap of \numCfgTrainingMaxEpochs{} passes was rerun with a cap of
\numCfgSensitivityEpochCapMaxEpochs{} passes, and the capped run was replaced by its twin (post hoc, rule logged
before the reruns). At \numCfgBudgetsHostsItemSix{} hosts $A$ moves from
\numSensCappedPostHocPerBudgetSixHundredFiftyFourPrimaryAMean{} to
\numSensCappedPostHocPerBudgetSixHundredFiftyFourCappedTwinsAMean{}~eV (paired change
\numSensCappedPostHocPerBudgetSixHundredFiftyFourChangeTwinMinusPrimaryAMean{}~eV
\ci{\numSensCappedPostHocPerBudgetSixHundredFiftyFourChangeTwinMinusPrimaryACiLo{}}{\numSensCappedPostHocPerBudgetSixHundredFiftyFourChangeTwinMinusPrimaryACiHi{}}),
and its upper bound from \numSensCappedPostHocPerBudgetSixHundredFiftyFourPrimaryAUpperOneSided{} to
\numSensCappedPostHocPerBudgetSixHundredFiftyFourCappedTwinsAUpperOneSided{}~eV, still below the tolerance. At
\numCfgBudgetsHostsItemFive{} hosts the bound stays above it under both versions
(\numSensCappedPostHocPerBudgetFourHundredPrimaryAUpperOneSided{} and
\numSensCappedPostHocPerBudgetFourHundredCappedTwinsAUpperOneSided{}~eV), so the smallest budget at which the
descriptors are bounded below the tolerance remains \numCfgBudgetsHostsItemSix{}. The twins cover only the runs of
the registered sweep; the capped runs of the extensions (Table~\ref{tab:capped}) were not rerun.

The architecture check trained an untuned CGCNN \citep{Xie2018crystal}, with fixed optimiser settings, on one outer
resample at three budgets (extension, descriptive). Its advantage falls with the number of training hosts as that of
both MEGNet networks does, from \numSensCgcnnTableCgcnnFiftyAPoint{}~eV at \numCfgBudgetsHostsItemTwo{} hosts to
\numSensCgcnnTableCgcnnSixHundredFiftyFourAPoint{}~eV
\ci{\numSensCgcnnTableCgcnnSixHundredFiftyFourACiNinetyFiveItemOne{}}{\numSensCgcnnTableCgcnnSixHundredFiftyFourACiNinetyFiveItemTwo{}}
at \numCfgBudgetsHostsItemSix{}, against \numSensCgcnnTableMegnetVOneSixHundredFiftyFourAPoint{}~eV for MEGNet v1
and \numSensCgcnnTableMegnetVTwoSixHundredFiftyFourAPoint{}~eV for MEGNet v2 on the same draws. CGCNN has
\numSensCgcnnParameterCountsCgcnnCgcnnSSixHundredFiftyFour{} trainable parameters in S, fewer than any of the tuned MEGNets
(\numSensCgcnnParameterCountsMegnetVTwoSTwoHundred{} to \numSensCgcnnParameterCountsMegnetVOneSTwoHundred{}). With one resample, the intervals reflect only the sampling of its test hosts, not split or retraining
variation, so the check shows that the fall of $A$ with training data is not peculiar to MEGNet; it does not
establish a bound for other architectures.

\begin{figure}[htbp]
\centering
\sifig{fig_si_sensitivities}
\caption{The advantage $A$ under the high-moment exclusion and the longer training cap, and an untuned CGCNN beside
the MEGNet networks. (\textbf{a}) $A = \mathrm{MAE}_\mathrm{S} - \mathrm{MAE}_\mathrm{D}$ (eV; positive when the DFT
descriptors lower the site error) at the largest training budget, from the seed-ensemble site MAE averaged over the
registered outer resamples. ``Primary'': the primary result of each network. ``High moment'' (extension): S and D
retrained with the entries whose defect magnetic moment exceeds the registered threshold removed from training and
test sites, compared on the common test sites. ``Capped twins'' (post hoc): every primary run that stopped at the
training cap replaced by its rerun with the longer cap. Both sensitivities exist for the originally registered
network only. Markers: point estimates; bars: two-sided intervals of the hierarchical bootstrap of the primary
analysis; short horizontal bars: one-sided upper bounds; dashed line: the tolerance $\delta$; dotted line: zero.
(\textbf{b}) Seed-ensemble test MAE (eV) of S (blue) and D (green) on a single outer resample at three budgets, for the
MEGNet networks (lines with bands) and for CGCNN with fixed, untuned optimiser settings (crosses, drawn just right of
their budget). Bands and bars in (b) are two-sided intervals of a host-cluster bootstrap over the test hosts of that
resample; with one resample they reflect test-host sampling only, and no test is made. Filled markers and solid
lines: primary network (v2), adopted after diagnostics of the originally registered network; open markers and
dashed lines: originally registered network (v1). Records: \texttt{results/sens\_high\_moment\_v1.json},
\texttt{results/sens\_capped\_post\_hoc.json}, \texttt{results/sens\_cgcnn.json}.}
\label{fig:sens}
\end{figure}
\FloatBarrier
